%------------------------------------------------------------------------------%

\usepackage{integracao_e_series}
\usepackage{cancel}

\title{Séries Numéricas -- Exemplos}

%------------------------------------------------------------------------------%
\begin{document}

\addtitle

%------------------------------------------------------------------------------%
%------------------------------------------------------------------------------%
\section{Séries Geométricas}

%------------------------------------------------------------------------------%
\begin{frame}{Séries Geométricas}
	Uma \emph{Série Geométrica} é uma série com a forma
	\[
    \alpha + \alpha r + \alpha r^2 + \alpha r^3 + \cdots + \alpha r^{n-1} + \cdots
    = \sum_{n=1}^\infty \alpha r^{n-1}
  \]

  \pause\qquad\(\alpha\in\R\)          \tabto{33mm} número fixo não nulo, $\alpha\neq 0$

	\vspace{0.5\baselineskip}
  \pause\qquad\(r\in\R\)               \tabto{33mm} \emph{razão} da série

  \vspace{0.5\baselineskip}
  \pause\qquad\(a_n = \alpha r^{n-1}\) \tabto{33mm} termo geral

	\pause\vspace{\baselineskip}
	Especial pois sabemos quando converge e qual sua soma
\end{frame}

%------------------------------------------------------------------------------%
\begin{frame}{Convergência da Série Geométrica -- Parte 1}
	Se $r=1$
  \pause
  temos
	\[
    S_n
    = \sum_{k=1}^{n} \alpha \,1^{k-1}
    \pause
    = \sum_{k=1}^{n} \alpha
    \pause
    = \underbrace{\alpha + \cdots + \alpha}_{n \text{ vezes}}
    \pause
    = n\alpha
  \]

	\pause
	\vspace{\baselineskip}
	Portanto $S_n$ diverge quando $n\to\infty$
\end{frame}

%------------------------------------------------------------------------------%
\begin{frame}{Convergência da Série Geométrica -- Parte 2}
	Se $r\neq 1$ temos
	\begin{align*}
		\onslide<+->{S_n      & = \alpha  +    \alpha r + \alpha r^2 + \cdots + \alpha r^{n-1}                \\[2mm]}
		\onslide<+->{rS_n     & = \hspace{7mm} \alpha r + \alpha r^2 + \cdots + \alpha r^{n-1} + \alpha r^{n} \\[2mm]}
		\onslide<+->{S_n-rS_n & = \alpha - \alpha r^{n} \\[2mm]}
		\onslide<+->{S_n(1-r) & = \alpha - \alpha r^{n}        }
	\end{align*}
	\onslide<+->{
    \[
      S_n = \frac{\alpha (1-r^n)}{1-r} \qquad\qquad (r\neq 1)
    \]
    }
\end{frame}

%------------------------------------------------------------------------------%
\begin{frame}{Calculando o Limite}
	Com $r\neq 1$, temos \(S_n = \frac{\alpha(1-r^n)}{1-r}\),
  \pause
  portanto
  \[
    \lim S_n
    \pause
    = \lim \frac{\alpha(1-r^n)}{1-r}
    \pause
    = \frac{\alpha}{1-r} -\frac{\alpha}{1-r} \lim r^n
  \]
  \pause
  \vspace{-\baselineskip}
	\begin{description}
		\item[$r=-1$] \tabto{10mm}
		\onslide<6->{$r^n=(-1)^n\;$ diverge quando $n\to\infty$ portanto $S_n$ diverge} \\[7mm]
		\item[$\abs{r}>1$] \tabto{10mm}
		\onslide<7->{$\abs{r^n}\to\infty$ quando $n\to\infty$ portanto $S_n$ diverge} \\[7mm]
		\item[$\abs{r}<1$] \tabto{10mm}
		\onslide<8->{$r^n\to 0$ quando $n\to\infty$ portanto \( S_n \to \frac{\alpha}{1-r} \)}
	\end{description}
\end{frame}

%------------------------------------------------------------------------------%
\begin{frame}{Séries Geométricas}
	Uma série geométrica
	\[
    \sum_{n=1}^\infty \alpha r^{n-1}
  \]
	é convergente se $\abs{r}<1$ com soma
	\[
    \sum_{n=1}^\infty \alpha r^{n-1} = \frac{\alpha }{1-r}
  \]
	e divergente se $\abs{r}\geq 1$
\end{frame}

%------------------------------------------------------------------------------%
%------------------------------------------------------------------------------%
\section{Séries Telescópicas}

%------------------------------------------------------------------------------%
\begin{frame}{Séries Telescópicas}
	Raramente conseguimos fórmulas exatas para a soma de séries

	\pause
	\vspace{\baselineskip}
	\emph{Séries Telescópicas} são um caso especial onde quase todos os termos se anulam
\end{frame}

%------------------------------------------------------------------------------%
\begin{frame}{Exemplo}
	Calcular a soma da série
	\[\sum_{n=1}^{\infty} \left( \frac{1}{n} - \frac{1}{n+1} \right) \]
\end{frame}

%------------------------------------------------------------------------------%
\begin{frame}{Somas Parciais}
	\begin{align*}
		\onslide<+->{S_n & = \sum_{k=1}^{n} \left( \frac{1}{k} - \frac{1}{k+1} \right)
			\\[3mm]}
		\onslide<+->{& = \left( \frac{1}{1} - \frac{1}{1+1} \right)
			+ \left( \frac{1}{2} - \frac{1}{2+1} \right)
			+ \cdots
			+ \left( \frac{1}{n-1} - \frac{1}{n} \right)
			+ \left( \frac{1}{n} - \frac{1}{n+1} \right)
			\\[3mm]}
		\onslide<+->{&
      = \frac{1}{1} - {\color<4->{blue}\frac{1}{2}}
			+ {\color<4->{blue}\frac{1}{2}} - {\color<5->{red}\frac{1}{3}}
      + {\color<5->{red}\frac{1}{3}}  - {\color<6->{gray}\frac{1}{4}}
			+ \cdots
      + {\color<6->{gray}\frac{1}{n-2}}  - {\color<7->{green}\frac{1}{n-1}}
			+ {\color<7->{green}\frac{1}{n-1}} - {\color<8->{magenta}\frac{1}{n}}
			+ {\color<8->{magenta}\frac{1}{n}} - \frac{1}{n+1}
			\\[3mm]}
		\onslide<9->{& = 1 - \frac{1}{n+1}}
	\end{align*}
\end{frame}

%------------------------------------------------------------------------------%
\begin{frame}{Calculando o Limite}
	\begin{align*}
		\onslide<+->{\sum_{n=1}^{\infty} \left( \frac{1}{n} - \frac{1}{n+1} \right)
                  & = \lim_{n\to\infty} S_n \\[3mm]}
		\onslide<+->{ & = \lim_{n\to\infty} \left( 1 - \frac{1}{n+1} \right)  \\[3mm]}
		\onslide<+->{ & = 1}
	\end{align*}
\end{frame}

%------------------------------------------------------------------------------%
\begin{frame}{Pode Não Ser Óbvio}
	Calcular a soma da série
	\[\sum_{n=1}^{\infty} \frac{1}{\;n(n+1)\;} \]
	\pause
	Usando frações parciais
	\[ \frac{1}{n(n+1)} = \frac{1}{n} - \frac{1}{n+1} \]
	\pause
	Temos
	\[
	\sum_{n=1}^{\infty} \frac{1}{\;n(n+1)\;}
	\;=\; \sum_{n=1}^{\infty} \left( \frac{1}{n} - \frac{1}{n+1} \right)
	\;=\; 1
	\]
\end{frame}

%------------------------------------------------------------------------------%
\section{Lista Mínima}

%------------------------------------------------------------------------------%
\begin{frame}{Lista Mínima}
  Estudar as Seção 6.1 da Apostila

  \vspace{\baselineskip}
  Exercícios: 8a-c, 10, 11a-b

  \vfill
  Atenção: A prova é baseada no livro, não nas apresentações
\end{frame}

%------------------------------------------------------------------------------%
\end{document}
%------------------------------------------------------------------------------%
